% Standalone includable research fragment. Requires amsmath, amssymb,
% amsthm and theorem/lemma/corollary environments in the parent document.
% This is a new result in the present audit, not an assertion of priority
% outside the supplied repository corpus.
\section{Finite aperiodicity without injectivity or induction}
\label{sec:finite-aperiodicity}

The scalar self-seeded draft charges successor injectivity as an
axiom.\cite{selfseed}
There is a useful reduction of its unresolved arithmetic problem: every
\emph{fixed finite} successor cycle is already excluded by its computation
schemas and Boolean equality, without using either injectivity or any
induction rule. This does not exclude a merger of two distinct predecessor
chains. The distinction is essential.

\subsection{Exact assumptions}

Work in intuitionistic first-order logic with actual equality and falsity.
The numerical language has $0,S,\epsilon$, and, for each fixed earlier-stage
term $h(v,\vec p)$, a scalar iteration symbol $F_h(n,\vec p)$ satisfying
\begin{align}
 F_h(0,\vec p)&=S(h(0,\vec p)),\label{eq:ac-F0}\\
 F_h(Sn,\vec p)&=h(F_h(n,\vec p),\vec p).
 \label{eq:ac-FS}
\end{align}
The other assumptions used in this section are precisely
\begin{gather}
 0=S0\longrightarrow\bot,\label{eq:ac-sep}\\
 \epsilon(x,y)=0\ \lor\ \epsilon(x,y)=S0,
 \label{eq:ac-range}\\
 \epsilon(x,y)=S0\quad\longleftrightarrow\quad x=y.
 \label{eq:ac-reflection}
\end{gather}
The ordinary logical rules are specified in the next paragraph; they are
part of this precursor, not omitted from its cost. Write $1=S0$ and
$\bar j=S^j0$ for an externally fixed nonnegative integer
$j$. There is no quantification over function descriptions. Each displayed
$h$ is a finite term, and its $F_h$ symbol is introduced at a later finite
stage. Numerical coverage, successor injectivity, predecessor, arbitrary
initialization, and induction are not assumptions of this section.

\paragraph{Ordinary logical rule inventory.}
In the precursor, formulas are generated from numerical equations and
$\bot$ by $\land,\lor,\to,\forall,\exists$. Contexts are finite reusable
assumption lists. The complete logical inventory is hypothesis membership;
conjunction introduction and its two projections; the two disjunction
injections and disjunction elimination with a common conclusion from its
two discharged branch assumptions; implication introduction by discharging
its antecedent and modus ponens; universal introduction at a parameter
fresh for the retained context and specialization at a term; existential
introduction with a term witness and elimination at a parameter fresh for
the retained context and conclusion; and inference of any formula from
$\bot$. Equality has reflexivity and capture-free replacement of an equal
numerical term in an arbitrary formula. Weakening, exchange, contraction,
renaming and cut are admissible structural rules. These ordinary rules
have no excluded-middle or double-negation-elimination axiom. Quantifier
substitutions avoid capture. The displayed computation, separation,
range and reflection formulas are the only nonlogical axioms used here;
the next section adds exactly contextual literal-equation induction.

\begin{lemma}[Zero is not a successor]
\label{lem:ac-zero}
For every $a$, the assumptions prove $Sa\ne0$.
\end{lemma}
\begin{proof}
For the constant-zero update put $Z(n)=F_{v.0}(n)$. The two computation
schemas give $Z(0)=1$ and $Z(Sa)=0$. An equality $Sa=0$ would therefore
give $0=1$, contradicting \eqref{eq:ac-sep} by symmetry. No range theorem
for $Z$ and no constructor coverage are used.
\end{proof}

\subsection{An explicit finite-table compiler}

For a constant update define
\[
 D(n,a):=F_{v.a}(n).
\]
The computation schemas give
\[
 D(0,a)=Sa,\qquad D(Sn,a)=a.
\]
For a Boolean term $b$, meaning that $b=0\lor b=1$ has been proved,
define the conditional increment
\[
 \operatorname{Inc}(b,a):=D(\epsilon(b,0),a).
\]
Reflection and separation prove $\epsilon(0,0)=1$ and
$\epsilon(1,0)=0$; in the latter case use the two-value range axiom and
rule out the alternative value $1$. Thus
\begin{equation}
 \operatorname{Inc}(0,a)=a,\qquad
 \operatorname{Inc}(1,a)=Sa.
 \label{eq:ac-inc}
\end{equation}
These equations do not cancel a successor.

Fix a positive ordinary integer $k$, and suppose that
$\bar0,\ldots,\bar k$ have already been proved pairwise distinct.
The following notation is an external finite expansion:
\[
 T_0(v):=0,\qquad
 T_{j+1}(v):=
 \operatorname{Inc}\bigl(\epsilon(v,\bar j),\,\cdot\,\bigr)^{j+1}
                      \bigl(T_j(v)\bigr)
 \quad(0\le j<k),\qquad
 h_k(v):=T_k(v).
 \tag{Finite cycle term}\label{eq:ac-cycle-term}
\]
Here the exponent $j+1$ means exactly $j+1$ written compositions of the
displayed unary term template, not an object-level recursion or a new
iteration operation. Consequently \eqref{eq:ac-cycle-term} is a canonical
finite term in $0,S,\epsilon,F$ after the abbreviations are expanded.

For a fixed $i<k$, substituting $v=\bar i$ makes its $i$th equality test
equal to $1$ and all other displayed tests equal to $0$. Exactly $i+1$
conditional increments are then performed. For $v=\bar k$ all tests are
zero. Hence the source theory proves the finite table
\begin{equation}
 h_k(\bar i)=\overline{i+1}\ (0\le i<k),\qquad
 h_k(\bar k)=0.
 \label{eq:ac-table}
\end{equation}
It also proves, for arbitrary $v$,
\begin{equation}
 h_k(v)=\bar0\ \lor\cdots\lor\ h_k(v)=\bar k.
 \label{eq:ac-h-range}
\end{equation}
Indeed, make the finitely many case distinctions supplied by
\eqref{eq:ac-range}. If some displayed test is $1$, reflection identifies
$v$ with its corresponding numeral; distinctness makes all the other
tests zero, and the preceding computation applies. If all tests are zero,
the output is zero. This is a finite use of intuitionistic disjunction
elimination. It assumes no excluded middle for quantified formulas.

This construction also describes the exact finite-table capacity used
here: a numeral-valued table is a finite sum of disjoint Boolean
indicators, with each unit contribution implemented by
\eqref{eq:ac-inc}. It does not provide a selector between arbitrary
numerical values or an object-level evaluator of variable-length tables.

\begin{lemma}[The finite update has no $k$-periodic point]
\label{lem:ac-finite-update}
Under the displayed numeral-distinctness assumption, the source proves
\[
                       h_k^{\,k}(w)\ne w
                       \qquad\text{for every }w.
\]
Again, the superscript denotes a fixed finite composition.
\end{lemma}
\begin{proof}
If $h_k^{\,k}(w)=w$, then $k\ge1$ and the range theorem
\eqref{eq:ac-h-range} imply that $w$ is one of
$\bar0,\ldots,\bar k$. On these $k+1$ values, the fixed table
\eqref{eq:ac-table} is one cycle. Starting at any of them and applying
that table $k$ times ends at a different displayed numeral. Each case
therefore contradicts the assumed distinctness. The modular description
of this finite table is explanatory: all cases are finite computations
in the source theory.
\end{proof}

\begin{theorem}[Every fixed successor cycle is impossible]
\label{thm:ac-aperiodicity}
For each fixed ordinary integer $k\ge1$, assumptions
\eqref{eq:ac-F0}--\eqref{eq:ac-reflection} prove
\[
                           \forall a\;S^k a\ne a.
\]
No instance of numerical induction and no successor-injectivity axiom is
used in this family of proofs.
\end{theorem}
\begin{proof}
We construct these finite proofs by external induction on $k$.
Suppose all positive cycle lengths less than $k$ have been excluded.
The numerals $\bar0,\ldots,\bar k$ are pairwise distinct: equality of
$\bar0$ with a positive numeral contradicts Lemma~\ref{lem:ac-zero};
and equality $\bar i=\bar j$ with $1\le i<j\le k$ would give
\[
                   S^{j-i}(\bar i)=\bar i,
\]
contradicting an earlier theorem because $1\le j-i<k$. For $k=1$ the
second case is empty, so the construction starts from separation alone.

Construct the term $h_k$ of \eqref{eq:ac-cycle-term}. Suppose
$S^k a=a$, and put $w=F_{h_k}(a)$. Exactly $k$ applications of the
successor computation schema, followed by equality congruence, give
\[
 w=F_{h_k}(a)=F_{h_k}(S^k a)=h_k^{\,k}(F_{h_k}(a))=h_k^{\,k}(w).
\]
This contradicts Lemma~\ref{lem:ac-finite-update}. Discharge the
assumption and universally generalize $a$.

The induction in the first sentence is in the metatheory and produces
one finite source proof for each fixed $k$. It is not a numerical
induction instance of the theory being simplified.
\end{proof}

\begin{corollary}
Every fixed pair of distinct numerals is provably distinct, and, for
fixed $0\le i<j$, the source proves $S^i a\ne S^j a$ for arbitrary $a$.
\end{corollary}
\begin{proof}
The numeral assertion is part of the construction. For the second
assertion use Theorem~\ref{thm:ac-aperiodicity} with the element $S^i a$
and cycle length $j-i$.
\end{proof}

\subsection{Precisely what remains open}

The theorem excludes cycles along one successor chain. It does not infer
$a=b$ from $Sa=Sb$ for two arbitrary elements. A directed successor
structure can have two different predecessor chains merge and still have
no finite cycle. This observation identifies the gap; it is not a
countermodel to the full scalar-iteration theory.

The finite-table construction has a different fixed term $h_k$ for each
external $k$. It does not supply a term for cycles of arbitrary internal
length. The latter would require a uniform arithmetic construction with
proved internal invariants, and could plausibly support a cancellation
argument. It cannot be obtained merely by quantifying over the finite
proofs above: the source has neither a uniform total evaluator nor an
arbitrary-initializer iterator among its assumptions.

In particular, the existing scalar self-seeded completion uses injectivity
when its capped trajectory proves
$S(V(m,S^3m))=S^3m$ and concludes $V(m,S^3m)=S^2m$. The aperiodicity
theorem does not justify this cancellation. Deleting injectivity from that
completion on the basis of the present theorem would therefore be
unproved. The result is a new independently useful consequence and a
more precise research target, not a completed smaller axiomatization.
